% !PRINTED 405
\weberpage{405}
%<<CONT

und in den beiden Ausnahmefällen $D=-4$, $D=-3$ ist der
letzte Ausdruck durch 2 oder durch 3 zu dividieren.

\gesperrt{Es durchläuft hierin $\mathfrak{m}$ die Gesamtheit der Ideale einer
Klasse (nach $[Q]$), und man sieht, daß dieser Grenzwert
 von dieser besonderen Klasse nicht abhängig ist.}

Es ist bisweilen zweckmäßig, in der Summe auf der linken
Seite der Formel (12) nicht bloß solche Ideale von $\mathfrak{m}$ auszuschließen
, die mit $Q$ einen Teiler gemein haben, sondern auch
die, die nicht teilerfremd zu $\varDelta$ sind. Die Formel (12) gilt unter
dieser Voraussetzung unverändert, wenn man unter den $r$ der
rechten Seite nicht nur die Primteiler von $Q$, sondern alle Primteiler
 von $D$ versteht.

Ist nämlich $\mathfrak{r}$ ein in $D$, aber nicht in $Q$ aufgehendes Primideal
, so ist $N(\mathfrak{r}) = r_1$ eine in $D$, aber nicht in $Q$ aufgehende
natürliche Primzahl, und es ist $N(\mathfrak{r}\mathfrak{m}) = r_1\,N(\mathfrak{m})$. Nimmt man
also von der Summe der linken Seite von $y$ die den $\mathfrak{r}\mathfrak{m}$ entsprechenden % [sic] Gedruckt ist „von y“; syntaktisch wäre ein Bezug auf die linke Summe zu erwarten.
 Glieder noch weg, so kommt rechts der Faktor

$\Bigl(1 - \frac{1}{r_1}\Bigr)$ hinzu.

5. \gesperrt{Der Grenzwert in der Formel (12) ist nur von der
Diskriminante $D$, nicht von der besonderen Klasse $M$
abhängig und ist stets von Null verschieden.}

\begin{center}
\S\,112. \gesperrt{Klassenzahl.}
\end{center}

Wir nehmen nun eine Funktion $F(z)$ an, die übrigens nur
für ganzzahlige Werte von $z$ definiert zu sein braucht, von der
wir voraussetzen, daß die unendliche Reihe

\weq{1}{\overset{\mathfrak{a}}{\Sigma} F(N\mathfrak{a}),}
in der $\mathfrak{a}$ die Gesamtheit der Ideale des Körpers $\Omega$ (oder nur
einen Teil davon) durchläuft, \gesperrt{unbedingt konvergent} sei. Wir
wollen, wenn es sich um Ordnungen $[Q]$ handelt, von $\mathfrak{a}$ alle die
\gesperrt{ausschließen, die nicht relativ prim zu $Q$ sind.}

Irgend eine natürliche zu $Q$ teilerfremde Zahl $m$ kommt
dann unter den $N(\mathfrak{a})$ nach § 93, 4.

\[ \overset{n}{\Sigma} (\varDelta, n) \]
mal vor, wenn $n$ die Teiler von $m$ durchläuft, und wir erhalten
die Formel:

\weq{2}{\sum F(N\mathfrak{a}) = \overset{m}{\Sigma} \overset{n}{\Sigma} (\varDelta, n)\,F(m),}
%CONT>>
